% =============================================================================
% Template: definition.tex
% Skill:    iacr-math-prose
%
% Use for: introducing a new primitive interface, a security notion, or a
%          hardness assumption. The statement is italicised by amsthm; the
%          body is plain.
%
% Invariants (do not delete):
%   - Every set, group, and field is typed at first use.
%   - The security parameter is the project-wide letter (lambda or n).
%   - The DST, if any hash is invoked, is explicit and ciphersuite-prefixed.
%   - The probability bound, if quoted, uses concrete inequalities (no `approx`).
%   - Cross-reference via \cref{def:<short-name>}.
%
% See:
%   - references/notation-discipline.md
%   - references/theorem-statement-style.md (for security-notion definitions)
% =============================================================================

\begin{definition}[\TODO{Short name of the notion, e.g. EUF-CMA security}]
\label{def:\TODO{short-label-here}}
A \TODO{primitive name, e.g. signature scheme}
$\TODO{\Sig} = (\TODO{\Setup}, \TODO{\Sign}, \TODO{\Verify})$
is \emph{\TODO{verbal name of the property, e.g. existentially unforgeable under
chosen-message attack}} if for every PPT adversary $A$,
\[
  \TODO{\Adv}^{\TODO{\textsf{NOTION}}}_{\TODO{\Sig}, A}(\secparam)
  \;\leq\; \TODO{\negl(\secparam)}.
\]
The advantage is defined as
$\TODO{\Adv}^{\TODO{\textsf{NOTION}}}_{\TODO{\Sig}, A}(\secparam)
 := \Pr[\TODO{\Expt}^{\TODO{\textsf{NOTION}}}_{\TODO{\Sig}, A}(\secparam) = 1]$,
where $\TODO{\Expt}$ proceeds as follows.
\begin{enumerate}
  \item \TODO{Step 1 of the experiment, e.g. $(\pk, \sk) \samples \Setup(1^\secparam)$.}
  \item \TODO{Step 2, e.g. $A$ is given $\pk$ and access to a signing oracle.}
  \item \TODO{Step 3, e.g. $A$ outputs $(m^*, \sigma^*)$.}
  \item \TODO{Final step, e.g. output $1$ iff $\Verify(\pk, m^*, \sigma^*) = 1$ and $m^* \notin Q$.}
\end{enumerate}
\end{definition}

% -----------------------------------------------------------------------------
% Variant: hardness assumption (no advantage definition; the assumption IS the
% advantage being negligible).
% -----------------------------------------------------------------------------

\begin{definition}[\TODO{q-SDH assumption / DDH assumption / LWE assumption}]
\label{def:\TODO{assumption-label}}
Let $\TODO{G}$ be \TODO{a cyclic group of prime order $p$ with generator $g$}.
The \emph{\TODO{decisional Diffie--Hellman (DDH)}} assumption in $\TODO{G}$
states that for every PPT distinguisher $D$,
\[
  \TODO{\Adv}^{\textsf{DDH}}_{\TODO{G}, D}(\secparam)
  := \left| \Pr[D(g, g^a, g^b, g^{ab}) = 1] - \Pr[D(g, g^a, g^b, g^c) = 1] \right|
  \;\leq\; \negl(\secparam),
\]
where the probabilities are over $a, b, c \samples \Z_p$ and the coins of $D$.

\TODO{Citation: e.g., Boneh, ``The Decision Diffie-Hellman Problem'',
ANTS 1998, Definition 1.}
\end{definition}
